\section{Related work}\label{sec:related-work}

The closest comparison is \citet[Section 3, Theorems 1--2 and equation (14)]{KennedyBalakrishnanRobinsWasserman2024}, who characterize supervised pointwise CATE estimation under H\"older smoothness. Write \(n\) for the labeled sample size, \leanref{sym:d}{\ensuremath{d}} for the covariate dimension, and \leanref{sym:alpha}{\ensuremath{\alpha}}, \leanref{sym:beta}{\ensuremath{\beta}}, and \leanref{sym:gamma}{\ensuremath{\gamma}} for the propensity, control-mean, and treatment-effect smoothness orders. With \(s_{\mathrm{pub}}=(\alpha+\beta)/2\), their absolute-error exponents are
\[
\frac{1}{1+d/(2\gamma)+d/(4s_{\mathrm{pub}})}
\quad\text{when}\quad
s_{\mathrm{pub}}<\frac{d/4}{1+d/(2\gamma)},
\qquad
\frac{\gamma}{2\gamma+d}
\quad\text{otherwise}.
\]
Published attainment requires their matrix regularity, smooth nuisance-estimation errors, boundedness conditions, and control of the third-order remainder involving covariate-distribution estimation. Their localized second-order projection correction balances effect approximation, nuisance approximation, and sampling variation.
% [Supervised pointwise minimax results](https://arxiv.org/pdf/2203.00837)

The algebraic comparison appears in \cref{obj:lem:published-cate-benchmark}. It identifies the supervised exponents obtained by specializing the sample-size expression to zero auxiliary records. Its interpretation follows the published source model, which permits a bounded covariate density and uses a construction-specific design in the lower-bound argument. The present \leanref{S-1}{potential-outcome law class with known uniform design}, specified in \cref{obj:def:primitive-class}, fixes the covariate distribution throughout. Within that class, \cref{obj:thm:marked-component-certificate,obj:lem:nuisance-frontier-converse} supply the original-record testing argument. These results establish the lower bound in the experiment actually observed by the decision maker, while the published benchmark supplies the numerical supervised comparison.

This distinction becomes consequential when information arrives in unequal amounts. The \leanref{S-2}{sampling experiment} combines \(n\) labeled records with \(m\) independent outcome-free covariate-treatment records; its sampling law is given in \cref{obj:ass:sample-law}. The pointwise absolute-error characterization in \cref{obj:thm:sharp-annotation-frontier} retains both record counts, and \cref{obj:thm:supplied-propensity} identifies the precision attainable when the propensity function is supplied. Together, these results connect the amount of treatment information to the accuracy of a conditional causal contrast under the paper's fixed design, overlap, and smoothness conditions.

Related semi-supervised treatment-effect analyses study population functionals through asymptotic expansions. \citet[Theorem 2.1 and Corollary 2.1]{ChakraborttyDai2022} separate labeled residual variation from averaging over the enlarged covariate sample for average treatment effect estimation. Their corollary addresses a vanishing labeled fraction with a correctly specified propensity model, retaining explicit nuisance remainders; asymptotic normality follows under the stated remainder conditions. The corresponding quantile treatment effect expansions additionally account for the density at the target quantile \citep[Theorem 3.1 and Corollary 3.1]{ChakraborttyDai2022}. These expansions explain robustness and efficiency for marginal causal functionals. The present risk characterization instead evaluates a conditional contrast at a fixed interior profile and quantifies worst-case absolute error over the specified H\"older class.
% [Semi-supervised ATE and QTE expansions](https://arxiv.org/pdf/2201.00468)

The information carried by auxiliary records also distinguishes the surrogate-assisted setting of \citet[Sections 2.1, 2.5--2.6, and 4]{ChengAnanthakrishnanCai2021}. Their records include outcome-predictive surrogate variables, potentially measured after treatment, alongside baseline covariates and treatment. An imputation-based ATE estimator exploits these variables, with robustness to imputation-model misspecification and local efficiency under their ideal semi-supervised model; perturbation resampling supports inference. Their electronic health records application illustrates the role of surrogate information when chart-reviewed outcomes are scarce. Here, auxiliary observations contribute covariate-treatment information, and their value is characterized through the pointwise risk and sample requirements in \cref{obj:thm:sharp-annotation-frontier,obj:prop:annotation-threshold}.
% [Surrogate-assisted ATE estimation](https://pmc.ncbi.nlm.nih.gov/articles/PMC7758040/)

A further comparison concerns the statistical criterion used to measure precision. \citet[Theorems 4.2 and 4.4 and Corollary 4.5]{Kato2026} derive two-sample ATE efficiency bounds for regular estimators and establish attainment through cross-fitted estimation under nuisance consistency and product-rate conditions. Their auxiliary sample contains covariates, and their analysis allows distinct labeled and auxiliary covariate distributions subject to support conditions. With coincident distributions and the stated evaluation weighting, the efficiency bound separates outcome residual variation from variation in the conditional effect across covariate profiles. Those conclusions characterize asymptotic variance at a fixed positive sampling fraction. The present minimax criterion, defined in \cref{obj:def:minimax}, measures pointwise absolute error across the primitive class and accommodates every auxiliary-record count in the stated experiment.
% [Regular-estimator ATE efficiency results](https://arxiv.org/html/2511.08303)

Supplied-predictor conditional estimation offers another route for using auxiliary information. \citet[Sections 2.4 and 3 and Theorems 1--3]{SuiZhouZhouDai2026} combine labeled observations, unlabeled covariates, and a supplied predictor to estimate conditional moment targets. Their Sobolev-equivalent reproducing kernel Hilbert space conditions support pointwise error bounds and confidence intervals, with prediction and residual components determining variance. Structure-adaptive CATE estimation uses different restrictions: \citet[Sections 2.3 and 6, Corollary 4]{Kim2026StructureAdaptiveCATE} give a pointwise squared-error result under bounded ambient-RKHS and overlap conditions when the contrast lies in the designated lower-complexity subspace; their broader analysis also treats integrated squared error, source conditions, known lower-dimensional representations, and model selection. Those results use supervised outcomes and RKHS/Sobolev structure. Our experiment instead fixes a H\"older class, measures pointwise absolute error, and studies the extra information supplied by outcome-free covariate-treatment records.
% [Supplied-predictor conditional inference](https://arxiv.org/html/2603.05575)
% [Structure-adaptive CATE estimation](https://arxiv.org/html/2602.18958)

Methodologically, the projection correction belongs to the higher-order influence-function tradition. Finite-rank orthogonal projections provide a tractable representation of second-order corrections in nonlinear functional estimation \citep[Lemma 3]{Robins2009}. Orthogonal regression supplies the conditional-effect estimating structure: residualized treatment and outcome regression underlie the R-learner \citep[Section 1]{Nie2021}, and local polynomial double-residual regression connects this structure to pointwise smoothness-based estimation \citep[Theorem 3]{Kennedy2023}. The \leanref{S-3}{rectangular projection estimator} in \cref{obj:def:moments,obj:def:estimator} combines localization with independent treatment and outcome roles. Its upper bound in \cref{obj:thm:rectangular-upper} makes the separate sampling contributions explicit, linking the projection architecture to the unequal information available from labeled and outcome-free records.
% [Finite-rank projection calculus](https://pmc.ncbi.nlm.nih.gov/articles/PMC3475538/)
% [Residualized conditional-effect regression](https://arxiv.org/pdf/1712.04912)
% [Local polynomial double-residual regression](https://arxiv.org/pdf/2004.14497)
